% Generated draft: requires mathematical and source review.
\documentclass[11pt]{article}
\usepackage[margin=25mm]{geometry}
\usepackage{fontspec}
\setmainfont{Latin Modern Roman}
\usepackage{amsmath,amssymb,mathtools,unicode-math}
\setmathfont{Latin Modern Math}
\usepackage{enumitem,needspace,xurl,hyperref}
\setlength{\parindent}{0pt}
\setlength{\parskip}{5pt}
\setlength{\emergencystretch}{3em}
\newcommand{\N}{\mathbb{N}}
\newcommand{\Z}{\mathbb{Z}}
\newcommand{\Q}{\mathbb{Q}}
\newcommand{\R}{\mathbb{R}}
\newcommand{\C}{\mathbb{C}}
\newcommand{\F}{\mathbb{F}}
\newcommand{\A}{\mathbb{A}}
\newcommand{\PP}{\mathbb{P}}
\newcommand{\E}{\mathbb{E}}
\newcommand{\eps}{\varepsilon}
\newcommand{\dd}{\,\mathrm{d}}
\begin{document}
\section*{UnsolvedMath ID: 30004142}
\textbf{Source problem code:} catalogue entry
\subsection{Definitions and mathematical statement}
The supplied formulation has an unresolved dimension inconsistency. A smooth quintic hypersurface \(Q\subset\mathbb P^4(\C)\) has complex dimension \(3\) and real dimension \(6\). Thus a real four-dimensional symplectic manifold cannot be diffeomorphic to \(Q\). Moreover, a Lagrangian submanifold of a real four-dimensional symplectic manifold has dimension \(2\), so its fundamental class cannot supply a generator in degree \(3\). Reading ``four-dimensional'' as complex dimension also conflicts with being a quintic hypersurface in \(\mathbb P^4\). The source title's reference to fourfolds does not resolve this inconsistency.

\medskip
The following is a conditional statement of the apparent homological target, requiring confirmation of the dimension correction and the geometric construction. Suppose that \((X,\omega)\) is a real six-dimensional symplectic manifold diffeomorphic to a smooth quintic hypersurface in \(\mathbb P^4(\C)\), and that \(\pi:X\to B\) is the intended Lagrangian fibration. On its regular locus, this means that the fibres are three-dimensional submanifolds on which \(\omega\) vanishes. The allowed singular fibres, the structure on \(B\) needed to define tropical hypersurfaces, and the meaning of a lift are not specified by the supplied record.

For these data, the proposed question is whether there exist a finite integer \(m\geq 1\), tropical hypersurfaces \(T_1,\ldots,T_m\) in \(B\), and smooth embeddings
\[
\iota_j:S^3\hookrightarrow X,\qquad \iota_j^*\omega=0\quad(1\leq j\leq m),
\]
such that each \(\iota_j(S^3)\) is obtained by lifting \(T_j\) in the intended, presently unspecified sense, and the homomorphism
\[
\Phi:\Z^m\longrightarrow H_3(X;\Z),\qquad
(a_1,\ldots,a_m)\longmapsto
\sum_{j=1}^{m}a_j(\iota_j)_*[S^3]
\]
is surjective. Here \([S^3]\in H_3(S^3;\Z)\) denotes an oriented fundamental class. A Lagrangian sphere means an embedded submanifold diffeomorphic to \(S^3\), rather than merely a rational homology sphere.

The record supplies \(\operatorname{rank}H_3(X;\Z)=204\). Its main statement asks for ``enough'' spheres, whereas its literature assessment refers to a requested family of exactly \(204\). Accordingly, whether the requirement is some finite \(m\) or specifically \(m=204\) remains unresolved. Rank alone does not establish that an integral homology group is free or that exactly that many elements generate it. The record also does not settle whether the question concerns every admissible pair \((X,\pi)\), a particular fibration, or the existence of a suitable fibration. No mutual disjointness condition on the generating spheres is stated. This conditional formulation does not resolve these qualifications.
\subsection{Short English statement}
The apparent question asks whether tropical hypersurfaces in the base of a Lagrangian fibration can be lifted to embedded Lagrangian \(3\)-spheres whose integral homology classes generate \(H_3(X;\Z)\). The supplied dimension is inconsistent with a smooth quintic in \(\mathbb P^4\); the six-dimensional interpretation, the lifting construction, the quantifiers, and the requirement of exactly \(204\) spheres need confirmation. The supplied partial-solution label is dated curation requiring independent review.
\subsection{Sources}
\begin{itemize}
\item Local UnsolvedMath registry, supplied complete record with permanent ID \(30004142\), source problem code \texttt{catalogue entry}, in the Oberwolfach Reports Open Problems collection. The supplied proposed year is \(2019\); no individual proposer is named.
\item Supplied report citation: \emph{Tropical Geometry: new directions} (2020), Oberwolfach Reports, pp.\ 1245--1308, DOI \texttt{10.4171/owr/2019/21}. \url{https://doi.org/10.4171/owr/2019/21}. This is the report-level provenance supplied by the registry; no precise problem page or passage is supplied.
\item Mak--Ruddat, supplied literature reference: \url{https://arxiv.org/abs/1904.11780}.
\item \emph{Forum of Mathematics Sigma} version, supplied literature reference: \url{https://doi.org/10.1017/fms.2020.54}.
\item Literature assessment dated exactly \texttt{2026-08-22}: supplied dated curation, not independently verified. The record labels the problem partially solved and reports that tropical constructions in the mirror quintic produce closed embedded Lagrangian rational homology spheres, including more than \(300\) mutually disjoint examples. It expressly says these results do not establish generation of \(H_3\) by the requested family of \(204\) Lagrangian spheres. The record requests independent verification of time-sensitive status and cited claims.
\item The supplied statement status is \texttt{corrected\_verified}, with a reported manual semantic, typographic, and LaTeX review. This is source metadata and does not resolve the mathematical inconsistencies identified here. No external source was consulted for this draft.
\end{itemize}
\end{document}
